% GENERATED FILE. Edit canonical lessons, then run npm run generate:books.
% canonical-source-hash: fnv1a64:31254d3ae0467c1e
% canonical-lessons: PA-C283-rah, PA-C283-sathi, PA-C283-pahunchna, PA-C283-phasna, PA-C283-hillna

\chapter{On the Road}
\label{ch:pa-a2-on-the-road}
\hlchaptermodality{283}

\begin{chapteropening}
\textbf{By the end of this chapter:} \emph{I can say when we will arrive.}

Complete the last lesson of chapter 283: I can say when we will arrive.
\end{chapteropening}

\section[\texorpdfstring{\pa{ਰਾਹ} (rāh)}{ਰਾਹ (rāh)}]{\pa{ਰਾਹ} (rāh) — a route; the way}

\label{lesson:PA-C283-rah}



\begin{quote}
Before the new one: say the Punjabi for to keep safe, then the Punjabi for to come out.
\end{quote}

\subsection*{You'll want to know: \pa{ਰਾਹ}}
\textbf{\pa{ਰਾਹ}} — \emph{rāh} — \textquotedblleft{}a route; the way\textquotedblright{}.

\begin{quote}
\pa{ਪਿੰਡ} \pa{ਦਾ} \pa{ਰਾਹ} \pa{ਲੰਮਾ} \pa{ਹੈ।}
\end{quote}

\begin{quote}
\emph{piṇḍ dā rāh lammā hai.}
\end{quote}

\begin{quote}
The way to the village is long.
\end{quote}

A masculine noun.

\begin{cousinweb}
\textbf{\pa{ਰਾਹ}} (\emph{rāh}) is a loan, usually traced to Persian \emph{rāh}, \textquotedblleft{}road, way\textquotedblright{}.
\end{cousinweb}

\begin{grammarlens}[title={What you've built}]
One more word for saying what comes next.
\end{grammarlens}

\subsection*{Your turn}
\begin{itemize}
\raggedright
  \item \emph{Say it:} \emph{rāh}
  \item \emph{Say it:} \emph{rāh}, once more
  \item \emph{Say it:} \emph{piṇḍ dā rāh lammā hai.} — The way to the village is long.
  \item \emph{Recall:} say the Punjabi for to keep safe, then the Punjabi for to come out, then say \emph{rāh} again
\end{itemize}

\begin{tcolorbox}[breakable,colback=teal!4,colframe=teal!35!black,title={Before you move on}]
What does \pa{ਰਾਹ} mean? (\textquotedblleft{}A route; the way\textquotedblright{}.) Say it once more.
\end{tcolorbox}
\section[\texorpdfstring{\pa{ਸਾਥੀ} (sāthī)}{ਸਾਥੀ (sāthī)}]{\pa{ਸਾਥੀ} (sāthī) — a companion; a comrade}

\label{lesson:PA-C283-sathi}



\begin{quote}
Before the new one: say the Punjabi for to come out, then the Punjabi for a route.
\end{quote}

\subsection*{You'll want to know: \pa{ਸਾਥੀ}}
\textbf{\pa{ਸਾਥੀ}} — \emph{sāthī} — \textquotedblleft{}a companion; a comrade\textquotedblright{}.

\begin{quote}
\pa{ਸਫ਼ਰ} \pa{ਵਿੱਚ} \pa{ਅਮਨ} \pa{ਮੇਰਾ} \pa{ਸਾਥੀ} \pa{ਹੋਵੇਗਾ।}
\end{quote}

\begin{quote}
\emph{safar vich Aman merā sāthī hovegā.}
\end{quote}

\begin{quote}
Aman will be my companion on the journey.
\end{quote}

A masculine noun.

\begin{cousinweb}
\textbf{\pa{ਸਾਥੀ}} (\emph{sāthī}) is inherited from a reconstructed \emph{*sārthin}, \textquotedblleft{}having part in a company\textquotedblright{}; CDIAL lists Punjabi \emph{sāthī}, \textquotedblleft{}comrade\textquotedblright{}, under it (CDIAL 13366).
\end{cousinweb}

\begin{grammarlens}[title={What you've built}]
One more word for saying what comes next.
\end{grammarlens}

\subsection*{Your turn}
\begin{itemize}
\raggedright
  \item \emph{Say it:} \emph{sāthī}
  \item \emph{Say it:} \emph{sāthī}, once more
  \item \emph{Say it:} \emph{safar vich Aman merā sāthī hovegā.} — Aman will be my companion on the journey.
  \item \emph{Recall:} say the Punjabi for to come out, then the Punjabi for a route, then say \emph{sāthī} again
  \item \emph{Write it:} copy \textbf{\pa{ਸਾਥੀ}} once, checking the whole word
\end{itemize}

\begin{tcolorbox}[breakable,colback=teal!4,colframe=teal!35!black,title={Before you move on}]
What does \pa{ਸਾਥੀ} mean? (\textquotedblleft{}A companion; a comrade\textquotedblright{}.) Say it once more.
\end{tcolorbox}
\section[\texorpdfstring{\pa{ਪਹੁੰਚਣਾ} (pahunchṇā)}{ਪਹੁੰਚਣਾ (pahunchṇā)}]{\pa{ਪਹੁੰਚਣਾ} (pahunchṇā) — to arrive; to reach}

\label{lesson:PA-C283-pahunchna}



\begin{quote}
Before the new one: say the Punjabi for a route, then the Punjabi for a companion.
\end{quote}

\subsection*{You'll want to know: \pa{ਪਹੁੰਚਣਾ}}
\textbf{\pa{ਪਹੁੰਚਣਾ}} — \emph{pahunchṇā} — \textquotedblleft{}to arrive; to reach\textquotedblright{}.

\begin{quote}
\pa{ਅਸੀਂ} \pa{ਸ਼ਾਮ} \pa{ਨੂੰ} \pa{ਪਿੰਡ} \pa{ਪਹੁੰਚਾਂਗੇ।}
\end{quote}

\begin{quote}
\emph{asīṁ shām nū̃ piṇḍ pahunchāṅge.}
\end{quote}

\begin{quote}
We will reach the village in the evening.
\end{quote}

Past (I …): \pa{ਮੈਂ} \pa{ਘਰ} \pa{ਪਹੁੰਚਿਆ}, \emph{maiṁ ghar pahunchiā}.

\begin{cousinweb}
\textbf{\pa{ਪਹੁੰਚਣਾ}} (\emph{pahunchṇā}) goes back to Sanskrit \emph{prabhūta}, \textquotedblleft{}much, great\textquotedblright{}; CDIAL lists Punjabi \emph{pahũcṇā}, \textquotedblleft{}to arrive\textquotedblright{}, under it (CDIAL 8716).
\end{cousinweb}

\begin{grammarlens}[title={What you've built}]
One more word for saying what comes next.
\end{grammarlens}

\subsection*{Your turn}
\begin{itemize}
\raggedright
  \item \emph{Say it:} \emph{pahunchṇā}
  \item \emph{Say it:} \emph{pahunchṇā}, once more
  \item \emph{Say it:} \emph{asīṁ shām nū̃ piṇḍ pahunchāṅge.} — We will reach the village in the evening.
  \item \emph{Recall:} say the Punjabi for a route, then the Punjabi for a companion, then say \emph{pahunchṇā} again
  \item \emph{Write it:} copy \textbf{\pa{ਪਹੁੰਚਣਾ}} once, checking the whole word
\end{itemize}

\begin{tcolorbox}[breakable,colback=teal!4,colframe=teal!35!black,title={Before you move on}]
What does \pa{ਪਹੁੰਚਣਾ} mean? (\textquotedblleft{}To arrive; to reach\textquotedblright{}.) Say it once more.
\end{tcolorbox}
\section[\texorpdfstring{\pa{ਫਸਣਾ} (phasṇā)}{ਫਸਣਾ (phasṇā)}]{\pa{ਫਸਣਾ} (phasṇā) — to get stuck; to be caught}

\label{lesson:PA-C283-phasna}



\begin{quote}
Before the new one: say the Punjabi for a companion, then the Punjabi for to arrive.
\end{quote}

\subsection*{You'll want to know: \pa{ਫਸਣਾ}}
\textbf{\pa{ਫਸਣਾ}} — \emph{phasṇā} — \textquotedblleft{}to get stuck; to be caught\textquotedblright{}.

\begin{quote}
\pa{ਬਰਸਾਤ} \pa{ਵਿੱਚ} \pa{ਗੱਡੀਆਂ} \pa{ਰਸਤੇ} \pa{ਵਿੱਚ} \pa{ਫਸਦੀਆਂ} \pa{ਹਨ।}
\end{quote}

\begin{quote}
\emph{barsāt vich gaḍḍīāṁ raste vich phasdīāṁ han.}
\end{quote}

\begin{quote}
In the rainy season, cars get stuck on the road.
\end{quote}

Past (I …): \pa{ਮੈਂ} \pa{ਰਸਤੇ} \pa{ਵਿੱਚ} \pa{ਫਸਿਆ}, \emph{maiṁ raste vich phasiā}.

\begin{cousinweb}
\textbf{\pa{ਫਸਣਾ}} (\emph{phasṇā}) is inherited from a reconstructed \emph{*spaśyatē}, \textquotedblleft{}is bound\textquotedblright{}; CDIAL lists Punjabi \emph{phasṇā}, \textquotedblleft{}to be ensnared\textquotedblright{}, under it (CDIAL 13812).
\end{cousinweb}

\begin{grammarlens}[title={What you've built}]
One more word for saying what comes next.
\end{grammarlens}

\subsection*{Your turn}
\begin{itemize}
\raggedright
  \item \emph{Say it:} \emph{phasṇā}
  \item \emph{Say it:} \emph{phasṇā}, once more
  \item \emph{Say it:} \emph{barsāt vich gaḍḍīāṁ raste vich phasdīāṁ han.} — In the rainy season, cars get stuck on the road.
  \item \emph{Recall:} say the Punjabi for a companion, then the Punjabi for to arrive, then say \emph{phasṇā} again
  \item \emph{Write it:} copy \textbf{\pa{ਫਸਣਾ}} once, checking the whole word
\end{itemize}

\begin{tcolorbox}[breakable,colback=teal!4,colframe=teal!35!black,title={Before you move on}]
What does \pa{ਫਸਣਾ} mean? (\textquotedblleft{}To get stuck; to be caught\textquotedblright{}.) Say it once more.
\end{tcolorbox}
\section[\texorpdfstring{\pa{ਹਿੱਲਣਾ} (hillṇā)}{ਹਿੱਲਣਾ (hillṇā)}]{\pa{ਹਿੱਲਣਾ} (hillṇā) — to shake; to sway}

\label{lesson:PA-C283-hillna}



\begin{quote}
Before the new one: say the Punjabi for to arrive, then the Punjabi for to get stuck.
\end{quote}

\subsection*{You'll want to know: \pa{ਹਿੱਲਣਾ}}
\textbf{\pa{ਹਿੱਲਣਾ}} — \emph{hillṇā} — \textquotedblleft{}to shake; to sway\textquotedblright{}.

\begin{quote}
\pa{ਪੁਰਾਣੀ} \pa{ਬੱਸ} \pa{ਬਹੁਤ} \pa{ਹਿੱਲਦੀ} \pa{ਹੈ।}
\end{quote}

\begin{quote}
\emph{purāṇī bass bahut hilldī hai.}
\end{quote}

\begin{quote}
The old bus sways a lot.
\end{quote}

Past (it …): \pa{ਉਹ} \pa{ਪੁਲ} \pa{ਹਿੱਲਿਆ}, \emph{oh pul hilliā}.

\begin{cousinweb}
\textbf{\pa{ਹਿੱਲਣਾ}} (\emph{hillṇā}) is inherited from a reconstructed \emph{*hillati}, \textquotedblleft{}moves, shakes\textquotedblright{}; CDIAL lists Punjabi \emph{hillṇā}, \textquotedblleft{}to shake, be moved\textquotedblright{}, under it (CDIAL 14120).
\end{cousinweb}

\begin{grammarlens}[title={What you've built}]
One more word for saying what comes next.
\end{grammarlens}

\subsection*{Your turn}
\begin{itemize}
\raggedright
  \item \emph{Say it:} \emph{hillṇā}
  \item \emph{Say it:} \emph{hillṇā}, once more
  \item \emph{Say it:} \emph{purāṇī bass bahut hilldī hai.} — The old bus sways a lot.
  \item \emph{Recall:} say the Punjabi for to arrive, then the Punjabi for to get stuck, then say \emph{hillṇā} again
  \item \emph{Write it:} copy \textbf{\pa{ਹਿੱਲਣਾ}} once, checking the whole word
\end{itemize}

\begin{tcolorbox}[breakable,colback=teal!4,colframe=teal!35!black,title={Before you move on}]
What does \pa{ਹਿੱਲਣਾ} mean? (\textquotedblleft{}To shake; to sway\textquotedblright{}.) Say it once more.
\end{tcolorbox}
